\section{Methodology}
% TODO: Insert project structure diagram.

To address the requirements outlined in \autoref{sec:requirement_analysis}, the project will proceed through the following stages.

\subsection{Formalisation of Proof Certificates in Rocq}

The first step is to define a formal representation of Marabou proof certificates in Rocq. This involves modelling all relevant components of the certificates, such as constraints, bounds, and the overall proof tree. The already formalised Marabou certificates in Imandra should serve as a foundation. However, since Rocq and Imandra are different languages, significant changes may be required.

An example of the formalisation of Marabou's proof tree is shown in \autoref{code:proof_tree}.

\begin{lstlisting}[language=Coq, caption={Marabou proof tree as an inductive data type in Rocq}, label={code:proof_tree}]
Require Import Reals.

Inductive Split : Type :=
  | Relu      : nat -> nat -> nat -> Split
  | SingleVar : nat -> R -> Split.

Inductive ProofTree : Type :=
  | Leaf : list R -> ProofTree
  | Node : Split -> ProofTree -> ProofTree -> ProofTree.
\end{lstlisting}


\subsection{Proof Certificates Parsing}

The system's entry point is parsing proof certificates produced by Marabou as JSON into Rocq data types. This step enables the transition from external certificate formats to an internal representation suitable for formal reasoning.
The parser will extract the following information:
\begin{itemize}
    \item \textbf{Upper bounds.} A representation of upper bounds associated with variables.
    \item \textbf{Lower bounds.} A representation of lower bounds associated with variables.
    \item \textbf{Tableau.} A representation of the linear constraint system maintained by Marabou.
    \item \textbf{Constraints.} A collection of constraints defining the verification problem.
    \item \textbf{Proof tree.} A structured representation of the proof steps leading to the final result.
\end{itemize}

\subsection{Proof Checker Algorithm Reconstruction}

Once the proof certificate has been serialised into Rocq's native data structures, the next step is to reconstruct the mechanism for validating it based on the Imandra implementation. In particular, the proof tree must be evaluated to determine whether it establishes satisfiability (SAT) or unsatisfiability (UNSAT).

The mentioned functionality involves implementing a function that recursively traverses the proof tree and verifies each step. For leaf nodes, the checker must ensure that a valid UNSAT witness can be constructed. For internal nodes, it must ensure that the inference steps are sound. This function corresponds to the \texttt{check\_tree} procedure described in \cite{certified_proof_checker}, adapted to the Rocq formalisation.

\subsection{Proof of Soundness}

The soundness of the proof checker is established by proving that if \texttt{check\_tree} returns UNSAT, then the verification property genuinely holds. To achieve this, the proof of soundness is structured around the DNN Farkas lemma and the soundness theorem itself.

\subsubsection{DNN Farkas Lemma Translation} \mbox{} \\
The DNN Farkas lemma, originally proven in Imandra by Desmartin et al. \cite{certified_proof_checker}, is the key result used to verify the leaves of the proof tree. It will be translated into Rocq by the following approach.

The proof will be carried out interactively in Rocq using tactic-based proving. Since Imandra relies heavily on automated reasoning, steps discharged automatically in Imandra must be decomposed into explicit proof steps in Rocq. To handle the real arithmetic reasoning involved, the \texttt{lra} tactic will be used to discharge linear arithmetic goals automatically where possible.

Any auxiliary lemmas that are not available in Rocq's standard library will be proven independently before tackling the main lemma, following a bottom-up proof development strategy.

